\section{Data and methods}
\label{sec:methods}
\paragraph{Data.} We use the \texttt{game} category of PopSign ASL v1.0 \citep{starner2023popsign}: 250 signs recorded on
smartphones, CC~BY~4.0, with a training split of \nSignersTrain{} signers, a validation split of \nSignersVal{} and a test split of \nSignersTest{}. The code checks that
no signer appears in two splits. Vocabularies of 10, 50, 100 and 250 signs are nested prefixes of a fixed random order of the
signs. A vocabulary is reported only if every one of its signs was processed; others are marked as not run. Videos are
downloaded, converted to landmarks and deleted.

\paragraph{Landmarks and features.} We use MediaPipe Holistic 0.10.14 \citep{lugaresi2019mediapipe} on frames resized to a
480 pixel short side. Coordinates are scaled by frame width and height before any angle is computed (the earlier
implementation used normalised coordinates, which distorts angles on portrait frames). A sequence is trimmed to the span
between the first and last frame with a detected hand. Feature set F1 is the 210 pairwise angles between the 21 hand
connection vectors per hand (identical information to the earlier 21$\times$21 angle matrix; Euclidean distances differ by a
constant factor). F2 is wrist-centred, scale-normalised hand shape. F3 combines F1 and F2, F4 adds hand location relative
to the shoulders, F5 adds upper-body pose and F6 adds face-contour landmarks. The main grid uses F1. Missing hands are
zero-filled in the main setting.

\paragraph{Methods.} \textbf{DTW 1-NN} compares a query with every stored reference using DTW with Euclidean local cost
and a Sakoe--Chiba band chosen on validation signers, and returns the label of the nearest reference. \textbf{Euclidean 1-NN}
compares sequences resampled to 32 frames. \textbf{DBA prototype} replaces each sign's references by their DTW barycentre (length equal to the median reference length).
The learned models are a 1D-CNN, a two-layer BiGRU and a three-layer Transformer encoder (\nParamsTrans{} to \nParamsGru{} parameters) trained
with cross-entropy for a fixed number of gradient steps with no early stopping, so no validation labels enter training. DTW
distances are computed by our own compiled implementation, which is tested against tslearn \citep{tavenard2020tslearn}.
Every method family chooses its own preprocessing on validation signers only (50 signs, 10 examples per sign, one draw), then it is fixed: how missing hands are handled (zero-fill, carry the last detection forward, or drop those frames), whether features are divided by their standard deviation over the references, and the DTW band (3, 5, 8 frames or none) for DTW; learning rate ($10^{-3}$ or $3\times10^{-4}$) and epoch budget (20 or 80) for the learned models. Training takes $\min(\max(\lceil \text{epochs}\cdot N/64\rceil,200),4000)$ gradient steps for $N$ training examples, so the largest conditions get only a few epochs.

\paragraph{Protocol.} References (or training examples) per sign $k\in\{1,2,5,10,20,\text{all}\}$ are drawn at random from the
training split, nested across $k$, with three draw seeds (which also seed network initialisation). Test queries are up to
three videos per test signer and sign, drawn once for the largest vocabulary and restricted to the signs of each smaller one, so all methods and vocabularies share them. A query on which no hand is ever detected is
counted as an error. We report top-1 accuracy with 95\% intervals from resampling the 8 test signers (10{,}000 resamples,
percentile intervals; cf. \citealp{field2007cluster}), and paired differences use the same resamples. With only 8 clusters
the intervals are wide and we do not report p-values. Two hypotheses were fixed in advance and judged at 97.5\%
(\S\ref{sec:results}); all other comparisons are exploratory.
